\lecture

\begin{guiding}
We want $H^*(G_n;\Z_2)$, together with generators we can compute with. The plan
has four movements: finish the cell structure and count cells; prove the
Leray--Hirsch theorem; use projective bundles to obtain a splitting principle;
and conclude that $H^*(G_n;\Z_2)$ is a polynomial ring, which finally
\emph{constructs} the Stiefel--Whitney classes and proves they are unique.
\end{guiding}

\section{Schubert cells, completed}

Recall the notation of Lecture 4: for a symbol
$\sigma=(\sigma_1<\dots<\sigma_n)$ we have $e(\sigma)\subset G_n(\R^m)$, the set
$e'(\sigma)$ of orthonormal frames with $x_i\in H^{\sigma_i}$, its closure
$\bar e'(\sigma)$ where $x_i\in\bar H^{\sigma_i}$, and
$d(\sigma)=\sum_i(\sigma_i-i)$.

We first record the tool that drives the induction.

\begin{lemma}\label{lem:rotation}
For unit vectors $u,v\in\R^m$ let $T(u,v)\in\Or(m)$ be the rotation carrying $u$
to $v$ and fixing $(\R u+\R v)^\perp$ pointwise. Then $T(u,v)(x)$ depends
continuously on $(u,v,x)$, and if $u,v\in\R^k$ then
$T(u,v)x\equiv x \bmod \R^k$ for every $x$.
\end{lemma}

\begin{proof}
On the plane $\R u+\R v$ the rotation is given by an explicit formula in $u,v$
(Gram--Schmidt in that plane followed by the planar rotation through the angle
between $u$ and $v$), and it is the identity on the orthogonal complement; both
pieces depend continuously on $(u,v)$. If $u,v\in\R^k$ then $T(u,v)$ is the
identity on $(\R^k)^\perp$ and preserves $\R^k$, so $T(u,v)x-x\in\R^k$.
\end{proof}

\begin{lemma}\label{lem:closed-cell}
$\bar e'(\sigma)$ is a closed cell of dimension $d(\sigma)$ with interior
$e'(\sigma)$.
\end{lemma}

\begin{proof}
\begin{strategy}
Induct on $n$. For $n=1$ the set is a closed hemisphere. For the inductive step,
observe that a frame in $\bar e'(\sigma_1,\dots,\sigma_{n+1})$ is a frame in
$\bar e'(\sigma_1,\dots,\sigma_n)$ together with one extra unit vector,
orthogonal to the previous ones and lying in $\bar H^{\sigma_{n+1}}$. Rotations
identify the space of possible extra vectors, which varies with the frame, with
one fixed closed cell; this exhibits the whole space as a product.
\end{strategy}

\step{1} \emph{Base case $n=1$.} Here
$\bar e'(\sigma_1)=\{x\in\R^{\sigma_1}\mid |x|=1,\ x_{\sigma_1}\ge0\}$, a closed
hemisphere of $S^{\sigma_1-1}$, hence a closed cell of dimension
$\sigma_1-1=d(\sigma_1)$, with interior the open hemisphere $e'(\sigma_1)$.

\step{2} \emph{The model for the new vector.} Fix
$\sigma_1<\dots<\sigma_{n+1}$ and let $b_i\in\R^m$ be the standard basis vector
with a $1$ in position $\sigma_i$, so that
$(b_1,\dots,b_n)\in e'(\sigma_1,\dots,\sigma_n)$. Put
\[
D=\big\{\,u\in\bar H^{\sigma_{n+1}}\ \big|\ |u|=1,\
b_1\cdot u=\dots=b_n\cdot u=0\,\big\}.
\]
The conditions say that $u$ lies in $\R^{\sigma_{n+1}}$, has last coordinate
$\ge0$, and has vanishing coordinates in the $n$ positions
$\sigma_1,\dots,\sigma_n$. So $D$ is a closed hemisphere of a sphere of
dimension $\sigma_{n+1}-n-1$, hence a closed cell of dimension
$\sigma_{n+1}-n-1$.

\step{3} \emph{Moving a frame to the standard one.} Given
$(x_1,\dots,x_n)\in\bar e'(\sigma_1,\dots,\sigma_n)$, set
\[
T_{(x)}=T(b_n,x_n)\circ\dots\circ T(b_1,x_1)\in\Or(m).
\]
Then $T_{(x)}b_i=x_i$ for every $i$. Indeed $T(b_j,x_j)$ with $j<i$ fixes
$b_i$, because $b_i$ is orthogonal to both $b_j$ and $x_j$: orthogonality to
$x_j$ holds since $x_j\in\R^{\sigma_j}$ has vanishing coordinate in position
$\sigma_i>\sigma_j$. Then $T(b_i,x_i)$ carries $b_i$ to $x_i$, and the later
rotations $T(b_l,x_l)$ with $l>i$ fix $x_i$ for the same reason. The map
$(x)\mapsto T_{(x)}$ is continuous by Lemma~\ref{lem:rotation}.

\step{4} \emph{The product decomposition.} Define
\[
\Psi\colon\bar e'(\sigma_1,\dots,\sigma_n)\times D
\longrightarrow \bar e'(\sigma_1,\dots,\sigma_{n+1})
\]
by
\[
\Psi\big((x_1,\dots,x_n),\,u\big)=\big(x_1,\dots,x_n,\ T_{(x)}u\big).
\]
This lands where claimed. The vector $T_{(x)}u$ is a unit vector; it is
orthogonal to each $x_i$, since
$x_i\cdot T_{(x)}u=T_{(x)}b_i\cdot T_{(x)}u=b_i\cdot u=0$ as $T_{(x)}$ is
orthogonal; and it lies in $\bar H^{\sigma_{n+1}}$, because by
Lemma~\ref{lem:rotation} each rotation used moves a vector only within
$\R^{\sigma_i}\subset\R^{\sigma_{n+1}}$, so
$T_{(x)}u\equiv u\bmod\R^{\sigma_{n+1}-1}$ and the $\sigma_{n+1}$-th coordinate
of $T_{(x)}u$ equals that of $u$, which is $\ge0$.

\step{5} \emph{$\Psi$ is a homeomorphism.} It is continuous by Step 3. It is
injective, since the first $n$ entries determine $T_{(x)}$ and hence
$u=T_{(x)}^{-1}(T_{(x)}u)$. It is surjective: given
$(x_1,\dots,x_{n+1})$ in the target, put $u=T_{(x)}^{-1}x_{n+1}$ and reverse the
computations of Step 4. Source and target are compact Hausdorff, so $\Psi$ is a
homeomorphism.

\step{6} \emph{Conclusion.} By induction $\bar e'(\sigma_1,\dots,\sigma_n)$ is a
closed cell of dimension $\sum_{i\le n}(\sigma_i-i)$; by Step 2, $D$ is a closed
cell of dimension $\sigma_{n+1}-(n+1)$. A product of closed cells is a closed
cell, and dimensions add, giving $d(\sigma_1,\dots,\sigma_{n+1})$. Under $\Psi$
the interiors correspond, since $T_{(x)}u$ has positive $\sigma_{n+1}$-th
coordinate exactly when $u$ does; hence the interior is
$e'(\sigma_1,\dots,\sigma_{n+1})$.
\end{proof}

\begin{lemma}\label{lem:q-characteristic}
The map $q|_{\bar e'(\sigma)}\colon\bar e'(\sigma)\to G_n(\R^m)$ has image
$\overline{e(\sigma)}$ and restricts to a homeomorphism
$e'(\sigma)\to e(\sigma)$. In particular it is a characteristic map for the cell
$e(\sigma)$.
\end{lemma}

\begin{proof}
By Claim~\ref{claim:half-space} the restriction $q\colon e'(\sigma)\to e(\sigma)$
is a bijection; it is continuous, and its inverse assigns to $X$ its unique
positive orthonormal basis, which is continuous by the explicit inductive
construction in the proof of that claim. Since $\bar e'(\sigma)$ is compact and
$G_n(\R^m)$ Hausdorff, $q(\bar e'(\sigma))$ is compact, hence closed, and it
contains $e(\sigma)$, so it contains $\overline{e(\sigma)}$; conversely
$\bar e'(\sigma)$ is the closure of $e'(\sigma)$, so its image lies in the
closure of $e(\sigma)$.
\end{proof}

\begin{theorem}\label{thm:cw-grassmannian}
The $\binom{m}{n}$ sets $e(\sigma)$ are the cells of a CW structure on
$G_n(\R^m)$, with $\dim e(\sigma)=d(\sigma)$. Passing to the limit $m\to\infty$
gives a CW structure on $G_n$ with one cell for each partition into at most $n$
parts.
\end{theorem}

\begin{proof}
\begin{strategy}
Everything except the boundary condition has been proved. For that condition one
shows that a plane in the closure of $e(\sigma)$ has a symbol $\tau$ which is
termwise smaller, hence a strictly smaller cell dimension unless $\tau=\sigma$.
\end{strategy}

\step{1} By Lemma~\ref{lem:closed-cell} and Lemma~\ref{lem:q-characteristic},
each $e(\sigma)$ is an open cell of dimension $d(\sigma)$ with characteristic
map $q|_{\bar e'(\sigma)}$, and the $e(\sigma)$ partition $G_n(\R^m)$.

\step{2} \emph{Boundaries drop dimension.} Let
$X\in\overline{e(\sigma)}\smallsetminus e(\sigma)$, say $X=q(x_1,\dots,x_n)$
with $(x_i)\in\bar e'(\sigma)\smallsetminus e'(\sigma)$. Since
$x_i\in\bar H^{\sigma_i}\subset\R^{\sigma_i}$, we get
$x_1,\dots,x_i\in X\cap\R^{\sigma_i}$ and therefore
\[
\dim\big(X\cap\R^{\sigma_i}\big)\ \ge\ i\qquad\text{for every } i .
\]
If $\tau$ is the symbol of $X$, this says exactly $\tau_i\le\sigma_i$ for all
$i$, whence
\[
d(\tau)=\sum_i(\tau_i-i)\ \le\ \sum_i(\sigma_i-i)=d(\sigma).
\]

\step{3} \emph{Strictness.} Equality would force $\tau_i=\sigma_i$ for all $i$,
i.e.\ $X\in e(\sigma)$, which is excluded. Hence
$\overline{e(\sigma)}\smallsetminus e(\sigma)$ is contained in the union of
cells of strictly smaller dimension, of which there are finitely many.

\step{4} \emph{Limit.} The inclusions $G_n(\R^m)\subset G_n(\R^{m+1})$ are
inclusions of subcomplexes, since a symbol for $\R^m$ is one for $\R^{m+1}$ and
the cells match. The limit therefore inherits a CW structure, with one cell for
each sequence $1\le\sigma_1<\dots<\sigma_n$, that is, for each partition
$0\le\sigma_1-1\le\dots\le\sigma_n-n$ into at most $n$ parts.
\end{proof}

\begin{corollary}\label{cor:HGn-rank}
For every $r$,
\[
\dim_{\Z_2}H^r\big(G_n;\Z_2\big)\ \le\
p_n(r):=\#\{\text{partitions of }r\text{ into at most }n\text{ parts}\}.
\]
\end{corollary}

\begin{proof}
The cellular cochain complex computing $H^*(G_n;\Z_2)$ has exactly $p_n(r)$
generators in degree $r$, by Theorem~\ref{thm:cw-grassmannian}; cohomology of a
complex is a subquotient of the complex.
\end{proof}

\begin{remark}[quoted]
In fact the mod $2$ cellular boundary operator vanishes identically, so the
inequality above is an equality. We shall not need this: the reverse inequality
will come for free from Theorem~\ref{thm:HGn} below. See Milnor--Stasheff,
\S 6, for the direct verification that all incidence numbers of Schubert cells
are even.
\end{remark}

\section{The Leray--Hirsch theorem}

\begin{theorem}[Leray--Hirsch]\label{thm:leray-hirsch}
Let $\pi\colon E\to B$ be a fibre bundle with fibre $F$ and
$i\colon F\hookrightarrow E$ the inclusion of a fibre. Assume that $B$ admits a
finite trivializing cover, that $H^*(F;\Lambda)$ is a free $\Lambda$-module of
finite rank, and that there is a $\Lambda$-linear map
\[
s\colon H^*(F;\Lambda)\longrightarrow H^*(E;\Lambda)
\qquad\text{with}\qquad
i^*\circ s=\mathrm{id}_{H^*(F)} .
\]
Then
\[
\Phi\colon H^*(F)\otimes_\Lambda H^*(B)\longrightarrow H^*(E),
\qquad
\alpha\otimes\beta\longmapsto s(\alpha)\smile\pi^*\beta,
\]
is an isomorphism. In particular $\pi^*\colon H^*(B)\to H^*(E)$ is injective.
\end{theorem}

\begin{proof}
\begin{strategy}
The map $\Phi$ is a homomorphism of $H^*(B)$-modules and is defined globally, so
it can be compared over the members of a cover. Over a trivializing open set
K\"unneth applies, and the matrix of $\Phi$ in the K\"unneth basis is triangular
with respect to the degree of the fibre classes, with invertible diagonal
blocks; hence $\Phi$ is an isomorphism there. Mayer--Vietoris and the five lemma
then propagate the statement from a cover to a union, and induction on the size
of the cover finishes.
\end{strategy}

\step{1} \emph{Restriction.} For open $U\subset B$ let $E_U=\pi^{-1}(U)$ and let
$\Phi_U$ be defined by the same formula, using the restrictions of $s(\alpha)$ to
$E_U$. Restriction of cohomology classes commutes with $\Phi$, so we obtain maps
of long exact sequences whenever we use Mayer--Vietoris.

\step{2} \emph{The trivial case.} Assume $E_U\iso U\times F$. Choose a
homogeneous basis $\alpha_1,\dots,\alpha_N$ of $H^*(F)$ ordered so that
$\deg\alpha_1\ge\dots\ge\deg\alpha_N$. By K\"unneth,
\[
H^*(E_U)\ \iso\ H^*(U)\otimes H^*(F),
\]
a free $H^*(U)$-module with basis $1\otimes\alpha_1,\dots,1\otimes\alpha_N$.
Write
\[
s(\alpha_j)\big|_{E_U}=\sum_{k}c_{jk}\otimes\alpha_k,
\qquad c_{jk}\in H^{\deg\alpha_j-\deg\alpha_k}(U).
\]
Since degrees are non-negative, $c_{jk}=0$ unless
$\deg\alpha_j\ge\deg\alpha_k$, so the matrix $(c_{jk})$ is block triangular for
our ordering. Restricting to a fibre kills all classes of positive degree
pulled back from $U$, so $i^*s=\mathrm{id}$ gives, for
$\deg\alpha_j=\deg\alpha_k$, that $c_{jk}\in H^0(U)$ has the value $\delta_{jk}$
at every point; hence the diagonal blocks are invertible over $H^0(U)$.

\step{3} \emph{Invertibility.} A block triangular matrix over a commutative
ring with invertible diagonal blocks is invertible, because its strictly
triangular part is nilpotent (it strictly decreases the finite ordering). Hence
$(c_{jk})$ is invertible over $H^*(U)$, and therefore
\[
\Phi_U(\alpha_j\otimes\beta)=s(\alpha_j)\smile\pi^*\beta
=\sum_k\big(c_{jk}\,\beta\big)\otimes\alpha_k
\]
is an isomorphism of $H^*(U)$-modules.

\step{4} \emph{Mayer--Vietoris.} Suppose $U=U'\cup U''$ with the theorem known
for $U',U''$ and $U'\cap U''$. Tensoring the Mayer--Vietoris sequence of
$(U',U'')$ with the free module $H^*(F)$ keeps it exact, and $\Phi$ maps it to
the Mayer--Vietoris sequence of $(E_{U'},E_{U''})$, commutatively by Step 1.
Three out of every four consecutive vertical maps are isomorphisms, so the five
lemma gives that $\Phi_U$ is an isomorphism.

\step{5} \emph{Induction.} Let $U_1,\dots,U_r$ be a finite trivializing cover of
$B$. The theorem holds over each $U_i$ by Step 2, hence over
$U_1\cup U_2$ by Step 4, since $U_1\cap U_2$ is trivializing as well. Assuming
it over $U_1\cup\dots\cup U_j$, apply Step 4 to
$U'=U_1\cup\dots\cup U_j$ and $U''=U_{j+1}$, whose intersection is
$\bigcup_{i\le j}(U_i\cap U_{j+1})$, a union of $j$ trivializing sets to which
the inductive hypothesis applies. After $r$ steps we reach $B$.

\step{6} \emph{Injectivity of $\pi^*$.} Taking $\alpha=1\in H^0(F)$ and noting
$s(1)=1$, since $i^*$ is injective in degree $0$ on the relevant summand, the
composite $\beta\mapsto\Phi(1\otimes\beta)=\pi^*\beta$ is the inclusion of a
direct summand, hence injective.
\end{proof}

\begin{remark}
For the infinite Grassmannians we apply the theorem to the compact
approximations $G_n(\R^{n+k})$, which do admit finite trivializing covers, and
pass to the limit; the cohomology of $G_n$ in each fixed degree agrees with that
of $G_n(\R^{n+k})$ for $k$ large, since the CW structures are compatible and no
new cells appear in low dimensions.
\end{remark}

\section{Projective bundles and the splitting principle}

\begin{definition}
For a real bundle $\pi\colon E\to B$ of rank $n$, the \emph{projective bundle}
$p\colon\mathbb P(E)\to B$ is the fibre bundle whose fibre over $b$ is the space
$\mathbb P(\pi^{-1}(b))\iso\RP^{n-1}$ of lines in the fibre.
\end{definition}

\begin{lemma}\label{lem:PE-injective}
For every rank $n$ bundle $E\to B$ over a base with a finite trivializing cover,
\[
H^*\big(\mathbb P(E);\Z_2\big)\ \iso\ H^*(B;\Z_2)\otimes H^*(\RP^{n-1};\Z_2)
\]
as $H^*(B)$-modules; in particular
$p^*\colon H^*(B;\Z_2)\to H^*(\mathbb P(E);\Z_2)$ is injective.
\end{lemma}

\begin{proof}
\begin{strategy}
Produce the splitting $s$ required by Leray--Hirsch from a canonical line bundle
living on $\mathbb P(E)$: its classifying map restricts on each fibre to the
inclusion $\RP^{n-1}\hookrightarrow\RP^\infty$, which is injective on
cohomology in the relevant range.
\end{strategy}

\step{1} \emph{The tautological line bundle.} Over the point
$(b,\ell)\in\mathbb P(E)$, with $\ell\subset\pi^{-1}(b)$ a line, take the line
$\ell$ itself. This defines a line bundle $\gamma\to\mathbb P(E)$; local
triviality is inherited from that of $E$ together with the local triviality of
$\gamma^1$ over $\RP^{n-1}$.

\step{2} \emph{Its classifying map.} Let
$f\colon\mathbb P(E)\to G_1=\RP^\infty$ classify $\gamma$. On a fibre
$F=\mathbb P(\pi^{-1}(b))\iso\RP^{n-1}$ the bundle $\gamma$ restricts to the
canonical line bundle of $\RP^{n-1}$, so $f|_F$ classifies it and is therefore
homotopic to the standard inclusion $j\colon\RP^{n-1}\hookrightarrow\RP^\infty$.

\step{3} \emph{The splitting.} With
$H^*(\RP^\infty;\Z_2)=\Z_2[a]$ and
$H^*(\RP^{n-1};\Z_2)=\Z_2[a]/(a^{n})$, the map $j^*$ sends $a^k$ to $a^k$ and is
an isomorphism in degrees $\le n-1$. Define
\[
s\colon H^*(\RP^{n-1};\Z_2)\to H^*(\mathbb P(E);\Z_2),
\qquad s(a^k)=f^*(a^k)\quad(0\le k\le n-1),
\]
extended linearly. By Step 2, $i_F^*s(a^k)=j^*(a^k)=a^k$, so
$i_F^*\circ s=\mathrm{id}$.

\step{4} Apply Theorem~\ref{thm:leray-hirsch}.
\end{proof}

\begin{definition}
A \emph{splitting map} for a rank $n$ bundle $E\to B$ is a map $f\colon B'\to B$
such that $f^*E\iso L_1\oplus\dots\oplus L_n$ with all $L_i$ line bundles, and
$f^*\colon H^*(B)\to H^*(B')$ injective.
\end{definition}

\begin{corollary}[Splitting principle]\label{cor:splitting}
Every real vector bundle admits a splitting map.
\end{corollary}

\begin{proof}
\begin{strategy}
Split off one line at a time on the projectivization, where the tautological
line bundle is by construction a sub-bundle of the pull-back; the rank drops by
one at each stage and injectivity is preserved because a composite of injections
is injective.
\end{strategy}

\step{1} On $\mathbb P(E)$ the tautological line bundle $\gamma$ of Step 1 above
is a sub-bundle of $p^*E$: over $(b,\ell)$ the fibre of $p^*E$ is
$\pi^{-1}(b)\supset\ell$. Choosing a Euclidean metric,
Theorem~\ref{thm:orthogonal-complement} gives
\[
p^*E\ \iso\ \gamma\oplus\gamma^{\perp},
\qquad \rk\gamma^\perp=n-1 .
\]

\step{2} By Lemma~\ref{lem:PE-injective}, $p^*$ is injective on cohomology.

\step{3} Repeat the construction with $\gamma^{\perp}$ in place of $E$: form
$\mathbb P(\gamma^\perp)$, split off a further line, and so on. After $n-1$
steps the pull-back of $E$ to
\[
B'=\mathbb P\big(\cdots\mathbb P(\gamma^\perp)\cdots\big)
\]
is a sum of $n$ line bundles, and the composite of the projections is injective
on cohomology, being a composite of injective maps.
\end{proof}

Any identity between characteristic classes which holds for sums of line bundles
therefore holds in general. This is the workhorse of the subject.

\section{Cohomology of the Grassmannian}

\begin{definition}
A polynomial $p\in\Z_2[a_1,\dots,a_n]$ is \emph{symmetric} if it is invariant
under all permutations of the variables. The \emph{elementary symmetric
functions} are
\[
\sigma_1=\sum_i a_i,
\qquad
\sigma_2=\sum_{i<j}a_ia_j,
\qquad\dots,\qquad
\sigma_n=a_1a_2\cdots a_n .
\]
\end{definition}

\begin{theorem}[fundamental theorem of symmetric polynomials, quoted]
The subring of symmetric polynomials in $\Z_2[a_1,\dots,a_n]$ is freely
generated by $\sigma_1,\dots,\sigma_n$:
\[
\Z_2[a_1,\dots,a_n]^{S_n}=\Z_2[\sigma_1,\dots,\sigma_n].
\]
\end{theorem}

Recall that $H^*(\RP^\infty;\Z_2)=\Z_2[a]$ and that, by K\"unneth,
\[
H^*\big((\RP^\infty)^n;\Z_2\big)=\Z_2[a_1,\dots,a_n]
\]
with each $a_i$ of degree $1$.

\begin{theorem}\label{thm:HGn}
Let $h_n\colon(\RP^\infty)^n\to G_n$ classify
$\gamma^1\oplus\dots\oplus\gamma^1$, one summand pulled back from each factor.
Then
\[
h_n^*\colon H^*(G_n;\Z_2)\longrightarrow\Z_2[a_1,\dots,a_n]
\]
is injective with image the ring of symmetric polynomials. Consequently
\[
H^*(G_n;\Z_2)=\Z_2\big[\tilde\sigma_1,\dots,\tilde\sigma_n\big],
\qquad
\tilde\sigma_i:=(h_n^*)^{-1}(\sigma_i)\in H^i(G_n;\Z_2),
\]
and $\dim_{\Z_2}H^r(G_n;\Z_2)=p_n(r)$ for every $r$.
\end{theorem}

\begin{proof}
\begin{strategy}
Injectivity comes from the splitting principle: a splitting map for the
universal bundle factors through $h_n$, and an injective map cannot factor
through a non-injective one. Symmetry of the image comes from permuting the
factors, which does not change the bundle being classified. Finally the image is
all of the symmetric polynomials by comparing dimensions with the cell count of
Corollary~\ref{cor:HGn-rank}.
\end{strategy}

\step{1} \emph{Injectivity.} Let $f\colon B'\to G_n$ be a splitting map for
$\gamma^n$ (Corollary~\ref{cor:splitting}), so
$f^*\gamma^n\iso L_1\oplus\dots\oplus L_n$ and $f^*$ is injective. Let
$g\colon B'\to(\RP^\infty)^n$ have $i$-th component a classifying map for
$L_i$; then
\[
g^*\big(\gamma^1\oplus\dots\oplus\gamma^1\big)\iso L_1\oplus\dots\oplus L_n
\iso f^*\gamma^n .
\]
So $f$ and $h_n\circ g$ both classify the same bundle over $B'$, hence are
homotopic by Theorem~\ref{thm:universal}, giving $f^*=g^*\circ h_n^*$. Since
$f^*$ is injective, so is $h_n^*$.

\step{2} \emph{The image is symmetric.} A permutation $\rho$ of the factors of
$(\RP^\infty)^n$ satisfies
$\rho^*(\gamma^1\oplus\dots\oplus\gamma^1)\iso\gamma^1\oplus\dots\oplus\gamma^1$,
since it merely reorders the summands. Hence $h_n\circ\rho\simeq h_n$ and
$\rho^*h_n^*=h_n^*$. As $\rho^*$ permutes $a_1,\dots,a_n$, every element of the
image is invariant under permutations, so
\[
\im h_n^*\ \subset\ \Z_2[\sigma_1,\dots,\sigma_n].
\]

\step{3} \emph{Dimension count.} In degree $r$, the symmetric polynomials have
$\Z_2$-dimension equal to the number of monomials
$\sigma_1^{d_1}\cdots\sigma_n^{d_n}$ with $\sum i\,d_i=r$, that is, to the
number $p_n(r)$ of partitions of $r$ into at most $n$ parts (take $d_i$ parts
equal to $i$). By Steps 1 and 2 together with
Corollary~\ref{cor:HGn-rank},
\[
\dim H^r(G_n;\Z_2)\ \le\ p_n(r)
\qquad\text{and}\qquad
\dim \im h_n^*=\dim H^r(G_n;\Z_2).
\]

\step{4} \emph{Surjectivity onto the symmetric part.} It remains to produce
classes mapping onto $\sigma_1,\dots,\sigma_n$. The mod-$2$ cellular boundary of
the Schubert complex vanishes (quoted above, Milnor--Stasheff \S 6), so equality
holds in Corollary~\ref{cor:HGn-rank}:
$\dim H^r(G_n;\Z_2)=p_n(r)$. Combined with Step 3 this forces
$\im h_n^*$ to be the whole degree-$r$ symmetric part, for every $r$. Hence
$h_n^*$ is an isomorphism onto $\Z_2[\sigma_1,\dots,\sigma_n]$ and we may define
$\tilde\sigma_i$ as its preimage of $\sigma_i$.
\end{proof}

\begin{remark}
Step 4 is the one place where we import a computation. It can be avoided: the
independent construction of Stiefel--Whitney classes by Steenrod squares in
Lecture 7 produces classes $w_i(\gamma^n)\in H^i(G_n;\Z_2)$ with
$h_n^*w_i(\gamma^n)=\sigma_i$, which gives the surjectivity of Step 4 directly,
and then the equality $\dim H^r=p_n(r)$ becomes a corollary rather than an
input. Either route closes the argument; only the order of exposition changes.
\end{remark}

\section{Existence and uniqueness of the SW classes}

\begin{theorem}
There is exactly one family of classes satisfying the axioms of
Definition~\ref{def:sw-axioms}, namely
\[
w_i(\xi)=f_\xi^*\big(\tilde\sigma_i\big),
\]
where $f_\xi\colon B\to G_n$ is a classifying map for $\xi$.
\end{theorem}

\begin{proof}
\begin{strategy}
Existence: three of the four axioms are immediate from the definition, and the
fourth --- the product formula --- is proved in Lecture 6 by computing the
effect on cohomology of the map $G_n\times G_{n'}\to G_{n+n'}$ classifying the
direct sum. Uniqueness: naturality reduces everything to the universal bundle,
the splitting principle reduces the universal bundle to sums of line bundles,
and the product formula reduces sums of line bundles to a single line bundle,
where normalization pins the answer down.
\end{strategy}

\step{1} \emph{Rank.} $\tilde\sigma_0=1$ and $H^*(G_n)$ has no generators beyond
$\tilde\sigma_n$, so $w_0=1$ and $w_i=0$ for $i>n$ by definition.

\step{2} \emph{Naturality.} If $\eta=g^*\xi$ then $f_\eta\simeq f_\xi\circ g$ by
Theorem~\ref{thm:universal}, so
$w_i(\eta)=f_\eta^*\tilde\sigma_i=g^*f_\xi^*\tilde\sigma_i=g^*w_i(\xi)$.

\step{3} \emph{Normalization.} The inclusion $\RP^1\hookrightarrow\RP^\infty=G_1$
classifies $\gamma^1_1$, and $\tilde\sigma_1=a$ generates
$H^1(\RP^\infty;\Z_2)$, whose restriction to $H^1(\RP^1;\Z_2)$ is the generator;
hence $w_1(\gamma^1_1)\neq0$.

\step{4} \emph{Product formula.} Proved in Theorem~\ref{thm:whitney} of the next
lecture.

\step{5} \emph{Uniqueness.} Let $w_i'$ be any family satisfying the axioms. By
naturality both families are determined by their values on the universal bundle
$\gamma^n$, so it suffices to prove $w_i(\gamma^n)=w_i'(\gamma^n)$. Since
$h_n^*$ is injective (Theorem~\ref{thm:HGn}), it suffices to compare after
applying $h_n^*$, that is, to compare the classes of
$\gamma^1\oplus\dots\oplus\gamma^1$ over $(\RP^\infty)^n$. By the product
formula, both are determined by the classes of a single $\gamma^1\to\RP^\infty$.
There only $w_1$ is in question, by the rank axiom, and naturality applied to
$\gamma^1_1\to\gamma^1$ together with the normalization axiom forces
$w_1'(\gamma^1)$ to be non-zero, hence equal to the generator $a$, which is
$w_1(\gamma^1)$.
\end{proof}
